\section{Literature Review}
\label{sec:lr}

We organise related work by the three research questions. Extraction-specific
claims are drawn from an adversarially-verified sweep (independent verifiers;
confidence noted in text where relevant); the graph and ideation clusters are
drawn from the reviewed corpus.

\subsection{RQ2 --- Text-based gap mining}
\label{sec:lr-rq2}

\paragraph{The task and its two shapes.} Work on ``gaps'' divides into
\emph{extraction} (locate the sentences stating a limitation or future-work
direction) and \emph{generation} (synthesise a limitations/future-work passage).
Extraction systems converge on a single architecture---\textbf{high-recall
structural/lexical retrieval $\rightarrow$ high-precision classifier
filter}---which is precisely the funnel this thesis uses. \citet{huwan2015}
pair a regex future-work extractor with a four-way classifier; \citet{zhangfws2022}
use two stages (a Naive-Bayes binary recogniser reaching Macro-F1 $90.7\%$ on
their corpus, then a SciBERT six-way typer at weighted-F1 $72.6\%$) and release a
labelled NLP/ACL future-work corpus (9{,}009 future-work vs 55{,}887 non-future-work
sentences). The closest analogue to our funnel is the randomized-controlled-trial
limitations system of \citet{rctlimits}: a keyword slice of candidate sentences
followed by a fine-tuned PubMedBERT filter reaching detection F1 $0.821$, scaled
to $\sim$12k articles with no per-paper LLM, and only marginally above a
rule-only baseline ($0.800$)---direct support for our design and a realistic
target. \citet{futuregen2025} state the pattern outright (regex high-recall
$\rightarrow$ LLM filter $\rightarrow$ RAG generation) and trim their input ``to
reduce API cost,'' i.e.\ they pay exactly the per-paper cost our cheap Stage~B is
designed to avoid.

\paragraph{Limitation datasets and generation.} \citet{limgen2024} harvest the
\emph{mandated} ``Limitations'' sections of 4{,}068 ACL papers as gold and
benchmark generative models---the mandated-section harvest is the
distant-supervision recipe we adopt (\S\ref{sec:method}). \citet{bagels} build a
limitations dataset over ACL/NeurIPS/PeerJ (verbatim author-stated plus
peer-review-supplemented) and benchmark generation by BERTScore/ROUGE and
coverage; its ACL split is the benchmark used in \S\ref{sec:results}. The same
group iterates on limitation \emph{generation}: topic-modelled aggregation
\citep{limtopic}, a per-paper graph with RAG \citep{generatingsuggesti2024}, and a
multimodal chart-limitation variant \citep{mitigatingvisual2024}. Sentence-level
classification of literature-review, problem, and method sentences is studied by
\citet{modellingclassifyi2025} and \citet{extractingproblem2026}; the latter
reports that fine-tuned models beat LLM in-context learning on this task and
diagnoses formulaic-expression shortcut-learning---a failure mode any cue-based
classifier inherits. \citet{gapmap2025} and \citet{canllms2025} push toward
\emph{implicit} gaps (not author-stated)---the harder, more useful frontier, and
the boundary of what a faithful \emph{extractor} (this thesis) deliberately does
not attempt. \citet{mixtureknowledge2024} bridge gap identification and graph
structure for review generation.

\paragraph{The hard problems, corroborated.} The literature independently
confirms our findings: gaps concentrate in \emph{terminal} sections
(Discussion/Future-Work/Conclusion) but some hide mid-paper; authors conflate
future-work with limitations, and $>$50\% of stated limitations cannot be
discerned from text alone; and while binary detection is easy, fine-grained
typing is hard for everyone (SciBERT wF1 $72.6\%$ on six future-work types;
PubMedBERT F1 $0.49$ on 24 limitation sub-types). Extraction datasets are
domain-locked (ACL, biomedical); none targets cross-domain arXiv AI/ML/math,
motivating our own small gold set.

\subsection{RQ1 --- Graph-based forecasting}
\label{sec:lr-rq1}

The reviewed RQ1 cluster frames ``finding promising directions'' as
graph/link-prediction-over-time, which motivates our gap-graph.

\paragraph{Knowledge-graph link prediction for trend forecasting.}
\citet{predictingfuture2022} (Science4Cast) forecast future concept links in a
100k-paper network; \citet{impact4cast2024} and \citet{forecastinghigh2024}
predict \emph{high-impact} future links on evolving knowledge graphs;
\citet{forecastingfuture2023} and the dual-link-prediction technology-opportunity
analysis of \citet{technologyopportun2023} do the same for AI and patents.
\citet{predictingnew2025} extract concept graphs with LLMs to predict new
materials-science directions.

\paragraph{Emerging-topic detection over dynamic graphs.}
\citet{trackingdynamics2020} track co-word-network dynamics for emerging topics;
\citet{studypredictabilit2025} predict new scholar topics from knowledge networks;
\citet{empiricalstudy} couple a knowledge graph with an LLM to reconstruct
scientific history and forecast trends.

\paragraph{Graph topic modelling and representation learning} supply the
machinery for our community step: graph-isomorphism and neural graph topic models
\citep{ginopic2024,ngtm2025,graphconvolutional2023}, graph-based topic-diversity
and multiview online clustering \citep{noveltopic2023,multiviewdeep2022}, and
self-supervised graph autoencoders and negative sampling
\citep{s2gae2023,understandingnegat2020}. \citet{linkprediction2024} and
\citet{literaturebased2025} tie graph structure directly to hypothesis
generation---the RQ1$\rightarrow$RQ3 bridge. Gap2Idea's contribution relative to
this cluster is to run the same graph/recombination machinery over
\emph{extracted research-gap nodes} (not concept or citation nodes) and to expose
\emph{bridge} (edge-betweenness) and \emph{frontier} structure as idea seeds.

\subsection{RQ3 --- LLM-driven ideation}
\label{sec:lr-rq3}

\paragraph{Knowledge-graph-grounded ideation.} Closest to our graph-seeded
generation with a human study is \citet{generationhuman2024}. \citet{chimera2025}
build a knowledge base of \emph{idea recombinations}; \citet{cam2023} mine
creative analogies; and \citet{ramonllull} realise combinatorial
theme$\times$domain$\times$method ideation---the explicit recombination
hypothesis.

\paragraph{Idea-development tools and grounding studies.} IdeaSynth
\citep{ideasynth2025}, SCI-IDEA \citep{sciidea2025}, workflow scaffolding
\citep{scaffoldingflexibl2025}, and CRISP-IM \citep{adaptingcrisp2021} evolve idea
facets with a human in the loop. Quality studies motivate our evidence-grounding
and human evaluation: \citet{canchatgpt2024} and \citet{comparingideation2024}
compare LLM and human ideation, and \citet{improvingresearch2025} report gains in
feasibility and quality from metadata grounding.

\paragraph{Surveys and role framings.} \citet{reviewllm2025} review 61
LLM-ideation studies; \citet{evolvingrole}, \citet{ai4research2025}, and
\citet{aiscience2025} frame LLMs as evaluator/collaborator/scientist and chart the
move toward agentic science. LLM-as-judge recurs across the cluster (including
ablation meta-evaluation, \citealp{abgen2025}); its standing caveat---%
self-evaluation bias when judge and generator share a model family---directly
motivates our cross-provider panel and human study.

\paragraph{Where Gap2Idea sits.} Individually, each component has precedent. The
contribution is their integration: a cost-efficient extraction funnel feeding a
gap-graph whose bridges and frontiers seed grounded, hallucination-gated,
multi-agent ideation, validated by a human-anchored evaluation
methodology---a deployed, reproducible pipeline rather than any single new model.
